\paragraph{Theorem24 (Approximate Counting in Quasi-Linear-Time Hierarchy).} There is a $\Sigma_5\mathsf{TIME}[n\cdot\operatorname{polylog}(n)]$ machine $M$ such that on input $(\varphi,K)$, accepts if $\varphi$ has at least $K$ satisfying assignments, and rejects if $\varphi$ has at most $K/2$ satisfying assignments.
\begin{proof}[Proof Sketch]
Using universal hashing [69], we can show that there is a $\mathsf{BPTIME}[n\cdot\operatorname{polylog}(n)]^{\mathsf{SAT}}$ machine $M'$ that approximates the number of satisfying assignments of $\varphi'$. Let $H$ be a family of pairwise-independent hash functions whose output range is $\{0,1\}^{\lceil\log K\rceil}$, that is computable in near-linear time. (See, e.g., [72, Problem3.3,(2)] for an example based on Toeplitz matrices.) The machine $M'$ randomly samples $h\sim H$ and $r\sim\{0,1\}^{\lceil\log K\rceil}$, and uses the $\mathsf{SAT}$ oracle to decide whether there is a satisfying assignment $x$ of $\varphi'$ such that $h(x)=r$.
Finally, the lemma follows from the fact that for every oracle $\mathcal O$, $\mathsf{BPTIME}^{\mathcal O}[t]\subseteq\Sigma_3\mathsf{TIME}[t\cdot\operatorname{polylog}(t)]^{\mathcal O}$ [74].
\end{proof}
